%-------------------------
% Resume in Latex
% Author : Jake Gutierrez
% Based off of: https://github.com/sb2nov/resume
% License : MIT
%------------------------

\documentclass[letterpaper,11pt]{article}

\usepackage{latexsym}
\usepackage[empty]{fullpage}
\usepackage{titlesec}
\usepackage{marvosym}
\usepackage[usenames,dvipsnames]{color}
\usepackage{verbatim}
\usepackage{enumitem}
\usepackage[hidelinks]{hyperref}
\usepackage{fancyhdr}
\usepackage[english]{babel}
\usepackage{tabularx}
\usepackage{fontawesome5}
\usepackage{multicol}
\setlength{\multicolsep}{-3.0pt}
\setlength{\columnsep}{-1pt}
\usepackage{amsmath}
\usepackage[dvipsnames]{xcolor}

\input{glyphtounicode}


%----------FONT OPTIONS----------
% sans-serif
% \usepackage[sfdefault]{FiraSans}
% \usepackage[sfdefault]{roboto}
% \usepackage[sfdefault]{noto-sans}
% \usepackage[default]{sourcesanspro}

% serif
% \usepackage{CormorantGaramond}
% \usepackage{charter}


\pagestyle{fancy}
\fancyhf{} % clear all header and footer fields
\fancyfoot{}
\renewcommand{\headrulewidth}{0pt}
\renewcommand{\footrulewidth}{0pt}

% Adjust margins
\addtolength{\oddsidemargin}{-0.6in}
\addtolength{\evensidemargin}{-0.5in}
\addtolength{\textwidth}{1.19in}
\addtolength{\topmargin}{-.7in}
\addtolength{\textheight}{1.4in}

\urlstyle{same}

\raggedbottom
\raggedright
\setlength{\tabcolsep}{0in}

% Sections formatting
\titleformat{\section}{
  \vspace{-4pt}\scshape\raggedright\large\bfseries
}{}{0em}{}[\color{black}\titlerule \vspace{-5pt}]

% Ensure that generate pdf is machine readable/ATS parsable
\pdfgentounicode=1

%-------------------------
% Custom commands
\newcommand{\resumeItem}[1]{
  \item\small{
    {#1 \vspace{-2pt}}
  }
}

\newcommand{\classesList}[4]{
    \item\small{
        {#1 #2 #3 #4 \vspace{-2pt}}
  }
}

\newcommand{\resumeSubheading}[4]{
  \vspace{-2pt}\item
    \begin{tabular*}{1.0\textwidth}[t]{l@{\extracolsep{\fill}}r}
      \textbf{#1} & \textbf{\small #2} \\
      \text{\small#3} & \text{\small #4} \\
    \end{tabular*}\vspace{-7pt}
}

\newcommand{\resumeSubheadingDouble}[6]{
  \vspace{-2pt}\item
    \begin{tabular*}{1.0\textwidth}[t]{@{}l@{\extracolsep{\fill}}r@{}}
      \textbf{\Large #1} & \textbf{\Large \small #2} \\
      \text{\small#3} & \text{\small #4} \\
      \textbf{#5} & \textbf{\small #6} \\
    \end{tabular*}\vspace{-7pt}
}


\newcommand{\resumeSubSubheading}[2]{
    \item
    \begin{tabular*}{0.97\textwidth}{l@{\extracolsep{\fill}}r}
      \textit{\small#1} & \textit{\small #2} \\
    \end{tabular*}\vspace{-7pt}
}

\newcommand{\resumeProjectHeading}[2]{
    \item
    \begin{tabular*}{1.001\textwidth}{l@{\extracolsep{\fill}}r}
      \small#1 & \textbf{\small #2}\\
    \end{tabular*}\vspace{-7pt}
}

\newcommand{\resumeSubItem}[1]{\resumeItem{#1}\vspace{-4pt}}

\renewcommand\labelitemi{$\vcenter{\hbox{\tiny$\bullet$}}$}
\renewcommand\labelitemii{$\vcenter{\hbox{\tiny$\bullet$}}$}

\newcommand{\resumeSubHeadingListStart}{\begin{itemize}[leftmargin=0.0in, label={}]}
\newcommand{\resumeSubHeadingListEnd}{\end{itemize}}
\newcommand{\resumeItemListStart}{\begin{itemize}}
\newcommand{\resumeItemListEnd}{\end{itemize}\vspace{-5pt}}

%-------------------------------------------
%%%%%%  RESUME STARTS HERE  %%%%%%%%%%%%%%%%%%%%%%%%%%%%


\begin{document}

%----------HEADING----------
% \begin{tabular*}{\textwidth}{l@{\extracolsep{\fill}}r}
%   \textbf{\href{http://sourabhbajaj.com/}{\Large Sourabh Bajaj}} & Email : \href{mailto:sourabh@sourabhbajaj.com}{sourabh@sourabhbajaj.com}\\
%   \href{http://sourabhbajaj.com/}{http://www.sourabhbajaj.com} & Mobile : +1-123-456-7890 \\
% \end{tabular*}

\begin{center}
    % {\Huge \scshape kodavanti rama sravanth} \\ 
    % {\Huge \scshape Rama Sravanth Kodavanti} \\ 
    {\Huge \scshape Sravanth Kodavanti} \\ 
    % \vspace{1pt}
    \bigskip
    Bengaluru, Karnataka, 560048, India  \\ \vspace{1pt}
    \bigskip
    
    \small \raisebox{-0.1\height}\faPhone\ +91-9550173710 ~ \href{mailto:ramasravanthkodavanti@gmail.com}{\raisebox{-0.2\height}\faEnvelope\  \underline{ramasravanthkodavanti@gmail.com}} ~ 
    \href{https://linkedin.com/in/sravanthk27/}{\raisebox{-0.2\height}\faLinkedin\ \underline{sravanthk27}}  ~
    % \href{https://github.com/Sravanth-k27}{\raisebox{-0.2\height}\faGithub\ \underline{github.com/Sravanth-k27}}
    \href{https://scholar.google.com/citations?user=SiOQawIAAAAJ&hl=en}{\raisebox{-0.2\height}\faGraduationCap\ \underline{Google Scholar}} ~
    \href{https://openreview.net/profile?id=~Sravanth_Kodavanti1}{\faPenNib \ \underline{OpenReview}} ~
    \href{https://sravanthkod.github.io/}{\raisebox{-0.2\height}\faGlobe\ \underline{Website}}

    \vspace{-8pt}
\end{center}
\vspace{-1pt}
\section{About}
% Research Engineer specializing in \textbf{post-training optimization} and \textbf{efficient AI} for real-world impact. Skilled in advanced \textbf{inference optimization}, \textbf{model compression} (quantization), and \textbf{Neural Architecture Search}, delivering performance improvements in both industrial and research projects. Passionate about bridging cutting-edge AI methods with scalable solutions to address practical business needs.
% Research Engineer specializing in \textbf{pre/post-training optimization} and \textbf{efficient AI} for real-world impact across \textbf{language} and \textbf{vision} models. Skilled in advanced \textbf{inference optimization}, \textbf{model compression} (quantization), and \textbf{Neural Architecture Search}, delivering performance improvements in both industrial and research projects. Passionate about bridging cutting-edge AI methods with scalable solutions to address practical business needs.

Research Engineer specializing in \textbf{AI model inference optimization on GPUs} and \textbf{efficient AI} for \textbf{language} and \textbf{vision} models. Experienced across the full efficiency stack --- from \textbf{GPU-scale inference optimization} to \textbf{on-device model compression} via \textbf{quantization}, \textbf{pruning}, \textbf{distillation}, and \textbf{Neural Architecture Search} --- delivering measurable speedups in both industrial deployments and research. Passionate about bridging cutting-edge AI methods with scalable, practical solutions.



% -----------EDUCATION-----------
\section{Education}
  \resumeSubHeadingListStart
  \resumeSubheading
        {Georgia Institute of Technology} {2026 - 2028}
      {Master of Science, \textbf{Major} - Computer Science, \textbf{Specilization} - Computational Perception and Robotics}{}
    \resumeSubheading
      {Indian Institute of Technology, Hyderabad}{2020 - 2024}
      {Bachelor of Technology, \textbf{Major} - Computer Science and Engineering, \textbf{Minor} - Entrepreneurship}{CGPA - 8.58}
  \resumeSubHeadingListEnd

\section{Research Publications}

\subsection*{Language:}

\begin{enumerate}
    \item \textit{\textbf{Sravanth Kodavanti}}, Sowmya Vajrala, Srinivas Miriyala \textit{et al}\\
    \href{https://arxiv.org/abs/2604.18655}{\textcolor{Blue}{\textbf{Unlocking the Edge Deployment and On-Device Acceleration of Multi-LoRA Enabled One-for-All Foundational LLM}}},
    \textit{ACL 2026}
\end{enumerate}

\subsection*{Vision:}

\begin{enumerate}

    \item Sowmya Vajrala, Akshay Bankar, Manjunath Arveti, Shreyas Pandith, \textit{\textbf{Sravanth Kodavanti}}, Subhajit Sanyal, Amit Unde, Srinivas Miriyala \\
    \href{https://arxiv.org/abs/2605.02767}{\textcolor{Blue}{\textbf{TOC-SR: Task-Optimal Compact diffusion for Image Super Resolution}}}, \href{https://lovif-eccv2026-workshop.github.io/}{\textit{LoViF@ECCV 2026}}

    \item Subhajit Sanyal*, Srinivas Miriyala*, Akshay Bankar*, Manjunath Arveti, Sowmya Vajrala, Shreyas Pandith, \textit{\textbf{Sravanth Kodavanti}}, Abhishek Ameta, Harshit, Amit Unde \\
    \href{https://arxiv.org/abs/2601.09823}{\textcolor{Blue}{\textbf{NanoSD: Edge Efficient Foundation Model for Real Time Image Restoration}}},\\
    \textit{CVPR 2026}

    \item \textit{\textbf{Sravanth Kodavanti}}, Manjunath Arveti, Sowmya Vajrala, Srinivas Miriyala, Vikram N R \\
    \href{https://arxiv.org/abs/2603.28405}{\textcolor{Blue}{\textbf{EdgeDiT: Hardware-Aware Diffusion Transformers for Efficient On-Device Image Generation}}}, \href{https://ai-benchmark.com/workshops/mai/2026/}{\textit{Mobile AI@CVPR 2026}}

    \item Sowmya Vajrala, \textit{\textbf{Sravanth Kodavanti}}, Srinivas Miriyala, \\
    \textcolor{Blue}{\textbf{Hardware-aware Low Light Image Enhancement on Edge}}, \href{https://www.wicv.org/}{\textit{WiCV@CVPR 2026}}

    \item Sowmya Vajrala, Akash Parmar, Prasanna R, \textit{\textbf{Sravanth Kodavanti}}, Manjunath Arveti, Srinivas Miriyala, Ashok Kumar \\ 
    \href{https://arxiv.org/abs/2603.29535}{\textcolor{Blue}{\textbf{Quantization with Unified Adaptive Distillation to enable multi-LoRA based one-for-all Generative Vision Models on edge}}},
    \href{https://ai-benchmark.com/workshops/mai/2026/}{\textit{Mobile AI@CVPR 2026}}

    \item Srinivas Miriyala*, Sowmya Vajrala*, Hitesh Kumar, \textit{\textbf{Sravanth Kodavanti}}, Vikram N R\\
    \href{https://arxiv.org/abs/2601.11684}
    {\textcolor{Blue}{\textbf{Mobile-friendly Image De-noising: Hardware-Conscious Optimization for Edge Application}}},\\
    \textit{ICASSP 2025}
    
\end{enumerate}

\subsection*{Preprints:}

\begin{enumerate}
        \item \textit{\textbf{Sravanth Kodavanti*}}, Manjunath Arveti \\
        \textcolor{Blue}{\textbf{MOSAIC: Multi-Objective Search for Allocating Inference Cache Across LLM Decoder Layers}}, \textit{Preprint}
        \item Sowmya Vajrala, \textit{\textbf{Sravanth Kodavanti}}, Manjunath Arveti, Shreyas Kasture, Utsav Tiwari, Uttam Kumar, Sriharsha, Vikram N R, Sharan Allur, Euntaik Lee, Dohyoung Kim, Rajib Basu \\
        \textcolor{Blue}{\textbf{Layer-Conditioned Forecast Token Modulation: A Lightweight Framework for On-Device Speculative Decoding}}, \textit{Preprint}
        \item \textit{\textbf{Sravanth Kodavanti}}, Srinivas Miriyala, Sowmya Vajrala, Vikram N R, Sharan Allur\\
    \href{https://arxiv.org/abs/2605.02794}{\textcolor{Blue}{\textbf{Edge-Efficient Image Restoration: Transformer Distillation into State-Space Models}}}, \\
    \textit{Preprint}
    
    \item Sowmya Vajrala*, Srinivas Miriyala*, \textit{\textbf{Sravanth Kodavanti}}\\
    \href{https://arxiv.org/abs/2601.11685}{\textcolor{Blue}{\textbf{Towards Efficient Image Deblurring for Edge Deployment}}},
    \textit{Preprint}
\end{enumerate}

% \resumeItemListEnd

\section{Patents}
% \resumeItemListStart
\begin{enumerate}
        \item \textit{\textbf{Sravanth Kodavanti*}}, Sowmya Vajrala, Srinivas Miriyala, Nikita Malik, Vikram N R\\
    \textcolor{Blue}{\textbf{(Mobile SWA) Multi - objective Optimization for Balanced Inference - time Latency - efficient and Eviction - aware SWA}}.\\
    \textbf{Provisional Specification Filed}. Proposes a dynamic window for generic token eviction algorithms reducing KV Cache without significant drop in accuracy of LLMs.
    \item \textit{\textbf{Sravanth Kodavanti*}}, Srinivas Miriyala*, Sowmya Vajrala*\\
    \textcolor{Blue}{\textbf{System and Method for Accelerating On-Device Large Language Model Inference}}.\\
    \textbf{Complete Specification Filed}. Proposes a novel self-speculative decoding technique to significantly improve efficiency and reduce latency in LLM inference on edge devices.
\end{enumerate}


% \resumeItemListEnd
%-----------EXPERIENCE-----------
\section{Work Experience}
  \resumeSubHeadingListStart


    \resumeSubheading
      {\href{https://nava.com/}{Nava}}{Oct 2026 -- Present}
      {\textbf{Research Engineer}}{}
    
    \resumeItemListStart
      \resumeItem{Working on \textbf{AI model inference optimization on GPUs} for large-scale models, spanning post-training optimization and efficient AI serving.}
    \resumeItemListEnd

        \resumeSubheading
      {Samsung Research Institute Bangalore}{Aug 2024 - Sep 2026}
      {\textbf{Machine Learning Research Engineer}}{\textit{\textbf{Mentor:} \href{https://in.linkedin.com/in/srinivas-s-miriyala}{\color{Blue}\textbf{Dr. Srinivas Miriyala}}}}
      \resumeItemListStart
      % \resumeItem{Working on \textbf{Samsung Neural Acceleration Platform} for AI Model acceleration \& their deployment in mobile devices. Using Techniques such as \textbf{Neural Architecture Search} (NAS) \& \textbf{Quantization} for model acceleration.}
      % \resumeItem{Using Techniques such as \textbf{Neural Architecture Search} (NAS) \& \textbf{Quantization} for model acceleration.}
      \resumeItem{Developing \textbf{Samsung Neural Acceleration Platform} for AI model acceleration and deployment on mobile devices, leveraging \textbf{Neural Architecture Search} (NAS) , \textbf{Quantization} and other techniques for performance optimization.}
      \resumeItem{\textbf{Commercialization \& Impact: }
      
      \begin{enumerate}
        \item \textbf{Accelerated inference speed} or the \textbf{Samsung LLM Gauss - L, 3B/3.5B Model} by implementing 
        \begin{itemize}

        \item \textbf{Speculative Decoding}, achieving a \textbf{2X improvement} in tokens per second (toks/sec), Successfully integrated into the \textbf{Samsung S25 series}. 
        \item Developed \textbf{PULSE}: a novel in-house \textbf{self-speculative decoding} algorithm, achieving a \textbf{5X improvement} in tokens per second (toks/sec). Method targeted for \textbf{Samsung S26 series}.
        \end{itemize}
          % \item \textbf{Accelerated on-device inference} for the \textbf{Samsung LLM Gauss - L, 3B Model} by implementing \textbf{Speculative Decoding}, achieving a \textbf{2X improvement} in tokens per second (toks/sec). Successfully integrated into the \textbf{Samsung S25 flagship series}.
        % \item \textbf{Developed NAS and quantization-based models} for \textbf{low-light video de-noising} (\textbf{2.5}× faster; deployed on \textbf{Samsung A56}), \textbf{demoire artifact removal} (\textbf{2.2}× faster; targeted for \textbf{Galaxy S26}) and \textbf{Text to Image: Stable Diffusion} (\textbf{2.3}× faster; targeted for \textbf{Galaxy S26}) delivering advanced, commercialization-ready imaging solutions.
          % \item  \textbf{Optimized low-light video de-noising model} using \textbf{NAS}, \textbf{Quantization}, achieving a \textbf{2.5× speedup in inference}. Successfully deployed in the \textbf{Samsung A56 device}.

    \item Developed \textbf{NAS and quantization} driven models delivering advanced, commercialization-ready imaging solutions, including:
    \begin{itemize}
        \item UNet-based model for \textbf{low-light video de-noising} (\textbf{2.5}× speedup; deployed on \textbf{Samsung Galaxy A56}),
        \item UNet-based model for \textbf{demoire artifact removal} (\textbf{2.2}× speedup; targeted for \textbf{Galaxy S26}),
        \item Stable Diffusion-based model for \textbf{text-to-image generation, motion photo enhancement, and related use-cases} (\textbf{2.3}× speedup; targeted for \textbf{Galaxy S26})
    \end{itemize}
      \end{enumerate} }
    % \resumeItem{\textbf{Research Publications:}
    % \begin{enumerate}
    %     \item \href{https://ieeexplore.ieee.org/document/10888855}{\color{Blue}\textbf{Mobile-friendly Image De-noising: Hardware Conscious Optimization for Edge Application}}, % \textbf{Mobile-friendly Image De-noising: Hardware Conscious Optimization for Edge Application}, 
    %     % \textit{IEEE ICASSP, } 2025.
    %     \textit{Proc. IEEE Int. Conf. Acoust., Speech, Signal Process. (ICASSP), }2025.
    % \end{enumerate}
    %     } 
% \resumeItem{\textbf{Patents:}
%     \begin{enumerate}
%         \item {\color{Blue}\textbf{System and Method for Accelerating On-Device Large Language Model Inference}}. \textbf{Provisional specification filed}. Proposes a novel self-speculative decoding technique to significantly optimize the efficiency and latency of LLM inference on edge devices.
%     \end{enumerate}
% }


\resumeItem{\textbf{Awards \& Recognition:}
% \resumeItem{\textbf{Awards \& Recognition:}
\begin{enumerate}

    % \item \textbf{Spot Award (Q4 2025)} – For contributions towards NAS commericilization \& Research Projects.
    
    \item \textbf{Advance Development(Q2 2026) \& Spot Award (Q2 2025)} – For novel speculative decoding algorithm enhancing LLM acceleration.
    
    \item \textbf{MD Project Incentive Award (2024–25)} – For AI model optimizations and deployment on edge devices.
    
    \item \textbf{Team Awesome Award (Q4 2024)} – For on-device optimization of \textbf{Samsung Gauss - L} in \textbf{S25} mobile series.
\end{enumerate}
}

% }

      % \resumeItem{Been part of AI Model Commercialization }
      \resumeItemListEnd
% \resumeSubheading
%   {Research Reviewer}{2025 -- Present}
%   {Peer-Reviewed Conferences and Workshops}{}
%   \resumeItemListStart
%     \resumeItem{Reviewed research submissions for \href{https://icml.cc/}{\color{Blue}\textbf{ICML}}, \href{https://neurips.cc/}{\color{Blue}\textbf{NeurIPS}},
%     \href{https://www.aclweb.org/portal/}{\color{Blue}\textbf{ACL}}
%     workshops \& conferences.}
%     \resumeItem{Provided thorough and constructive feedback to help ensure high standards of academic quality and rigor.}
%   \resumeItemListEnd



        \resumeSubheading   
      {Stealth Startup - \href{https://gikagraph.ai/}{\color{Blue}\textbf{Gika Graph.AI}}}{June 2024 - Aug 2024}
      {\textbf{AI Researcher}}{\textit{\textbf{Mentor:} \href{https://www.linkedin.com/in/manoj-agarwal-48b12a1/}{\color{Blue}\textbf{Dr. Manoj Agarwal}}}}
      \resumeItemListStart
      % \resumeItem{Worked on a thesis for proving \textbf{Knowledge Graphs} (KG) are better than \textbf{Vector Databases} in \textbf{Retrieval Augmented Generation} (RAG).}
      \resumeItem{Developed domain-specific search engines leveraging \textbf{Knowledge Graphs} (KG) to enhance semantic accuracy. Improved \textbf{Retrieval Augmented Generation} by using KG over vector databases for more precise, context-aware retrieval.}
      \resumeItem{Handled \textbf{Coreference} and \textbf{Entity Resolution} for fine-tuning data using LLMs such as \textbf{GPT-4, GPT-4o}, \textbf{LLAMA3}, \textbf{SpanBERT}, and \textbf{LingMess} with the \textbf{SpaCy} AI framework.}
      \resumeItemListEnd
       \vspace{5pt}
      \resumeSubheading
      {Hexagon R \& D India}{Jan 2024 - June 2024}
      {\textbf{Machine Learning Intern}}{\textit{\textbf{Mentor:} \href{https://in.linkedin.com/in/ankan-sengupta-a599b910}{\color{Blue}\textbf{Ankan Sengupta}}}}
      \resumeItemListStart
      \resumeItem{Developed a website for implementing segmentation $\&$ classification based on tags of various manufacturing plant sketches.Used various segmentation algorithms such as \textbf{RANSAC}, \textbf{DBSCAN}, \textbf{K-Means}.}
      \resumeItem{Used \textbf{Azure form recognizer} models \& other OCR models for tag identification.}
      \resumeItem{Developed a \textbf{flask website} for text recognition on manufacturing plant sketches using \textbf{Form, Doc Recognizer} \& also used OCR models for tag identification.}
      \resumeItem{ Involved in an \textbf{LLM }research project. Compared the results for text generation between \textbf{Mixtral}, \textbf{Mistral}, \textbf{LLAMA2}.}
      \resumeItemListEnd
% `    \vspace{5pt}
% \\
    \resumeSubheading
      {OnePlus / Oppo (OPLUS) Mobiles India R \& D}{Jan 2023 - June 2023}
      % {\textbf{Deep Learning Intern - Device AI}}{}
      {\textbf{Research Intern - Device AI}}{\textit{\textbf{Mentor:} \href{https://www.linkedin.com/in/dr-c-shyam-anand-a9252039/}{\color{Blue}\textbf{Dr. C Shyam Anand}}}}
      \resumeItemListStart
        \resumeItem{Implemented model compression techniques called \textbf{Quantization} , \textbf{Pruning }\& \textbf{Distillation} on various deep learning models such as \textbf{ResNet} , \textbf{Yolo} , \textbf{ViTs} , \textbf{ConvNeXt} , \textbf{Stable Diffusion Models} \& \textbf{Large Language Models} (LLAMA) . \textbf{Compressed all these models for deploying in mobile devices.}}
        \resumeItem{ Involved in research work on \textbf{Neural Style Transfer} (NST).}
        \resumeItem{Impemented a DL model for \textbf{LaTeX - OCR} task . The model's aim is to detect \& recognize the mathematical equations present in a research paper. Compressed the model for the integration with edge devices.}
      \resumeItemListEnd
         \vspace{5pt}
    % \resumeSubheading
    %   {Department of Computer Science , IIT Hyderabad}{Nov 2021 - Apr 2024}
    %   {\textbf{Teaching Assistant}}{}
    %   \resumeItemListStart
    %      \resumeItem{I have worked as a \textbf{TA} for the courses \textbf{Operating Systems} (CS3510) under \textit{Prof.} Sathya Peri , \textbf{Discrete Maths} (CS1010) under \textit{Prof.} Rakesh Venkat , \textbf{DBMS} (CS3550) under \textit{Prof.} Manish Singh \& \textbf{Theory of Computation} (CS2030) under \textit{Prof.} Subramanyam Kalyanasundaram. }
    % \resumeItemListEnd
 
    
  \resumeSubHeadingListEnd
% \vspace{-10pt}
\section{Teaching Experience}
    \resumeSubHeadingListStart
    \resumeSubheading
      {Department of Computer Science , IIT Hyderabad}{Nov 2021 - Apr 2024}
      {\textbf{Teaching Assistant}}{}
      \resumeItemListStart
         % \resumeItem{I have worked as a \textbf{TA} for the courses \textbf{Operating Systems} (CS3510) under \textit{Prof.} Sathya Peri , \textbf{Discrete Maths} (CS1010) under \textit{Prof.} Rakesh Venkat , \textbf{DBMS} (CS3550) under \textit{Prof.} Manish Singh \& \textbf{Theory of Computation} (CS2030) under \textit{Prof.} Subramanyam Kalyanasundaram. }
         \resumeItem{Served as a \textbf{Teaching Assistant} for multiple courses, including:
\begin{enumerate}
\item \textbf{Operating Systems} (CS3510) under \textit{\href{https://people.iith.ac.in/sathya_p/}{\color{Blue}\textbf{Dr. Sathya Peri}}}
\item \textbf{Discrete Mathematics} (CS1010) under \textit{\href{https://people.iith.ac.in/rakeshvenkat/}{\color{Blue}\textbf{Dr. Rakesh Venkat}}}
\item \textbf{Database Management Systems} (CS3550) under \textit{\href{https://people.iith.ac.in/msingh/}{\color{Blue}\textbf{Dr. Manish Singh}}}
\item \textbf{Theory of Computation} (CS2030) under \textit{\href{https://people.iith.ac.in/subruk/}{\color{Blue}\textbf{Dr. Subrahmanyam Kalyanasundaram}}}
\end{enumerate}
}
    \resumeItemListEnd
    \resumeSubHeadingListEnd
    \vspace{-10pt}
\section{Research Community Service}

\resumeSubHeadingListStart
\resumeSubheading
  {Research Reviewer}{2025 -- Present}
  % {Peer-Reviewed Conferences and Workshops}{}
  {}{}
  \resumeItemListStart
    \resumeItem{\textbf{Reviewer}:  
    \href{https://icml.cc/}{\color{Blue}\textbf{ICML}}, 
    \href{https://neurips.cc/}{\color{Blue}\textbf{NeurIPS}},
    \href{https://iclr.cc/}{\color{Blue}\textbf{ICLR}},
    \href{https://www.aclweb.org/portal/}{\color{Blue}\textbf{ACL}},
    \href{https://aaai.org/conference/aaai/}{\color{Blue}\textbf{AAAI}} \&
    \href{https://2026.ieeeicassp.org/}{\color{Blue}\textbf{ICASSP}}
    workshops and main tracks.}
    % \resumeItem{Provided constructive feedback to maintain academic rigor and contribute to the research community.}
    \resumeItem{\textbf{Area Chair}: \href{https://2026.ieeeicassp.org/}{\color{Blue}\textbf{ICASSP'26}}}
  \resumeItemListEnd
  \resumeSubHeadingListEnd
% \section{Research Community Service}
% \resumeSubheading
% {Research Reviewer}{2025 -- Present}
% {}{}
% \vspace{-10pt}
% \resumeItemListStart
% \resumeItem{Reviewed submissions for \href{https://icml.cc/}{\color{Blue}\textbf{ICML}}, \href{https://neurips.cc/}{\color{Blue}\textbf{NeurIPS}}, and \href{https://www.aclweb.org/portal/}{\color{Blue}\textbf{ACL}} workshops.}
% \resumeItem{Provided constructive feedback to ensure quality and rigor.}
% \resumeItemListEnd


%-----------PROJECTS-----------
\section{Projects}
%     \vspace{-5pt}
    \resumeSubHeadingListStart
          \resumeProjectHeading
      {\textbf{Continual Learning for 3D Point Cloud}}{\textit{\textbf{Guide:} \href{https://sites.google.com/site/pksrijith/home}{\color{Blue}\textbf{Dr. P K Srijith}}}}
      \resumeItemListStart
        \resumeItem{Implemented Continual Learning on \textbf{Pointnet} architecture by addressing the challenge of catastrophic forgetting.}
        \resumeItem{Implemented \textbf{Knowledge Distillation} approach for the Continual Learning \& Model was trained on \textbf{ModelNet10} dataset.}
          \resumeItemListEnd
          \vspace{-15pt}
      \resumeProjectHeading 
        {\textbf{Computer Vision \& NLP} $|$ Personal Projects}{}
        \resumeItemListStart
         \resumeItem{Image denoising using \textbf{Auto - Encoders }. Model trained on \textbf{MNIST} dataset.}
         \resumeItem{Human Face generation using \textbf{GANs} . Model trained on \textbf{celeba} dataset.}
         \resumeItem{Underwater Object Detection and Classification using \textbf{DGYOLO}. Model trained on \textbf{URPC 2019} dataset.}
          \resumeItemListEnd
%            \vspace{-15pt}
%         \newpage
%         \resumeProjectHeading
%        {\textbf{Optimizing B - Spline curves for robot navigation - IoT} $|$\textit{ Prof.} P.Rajyalakshmi} {}
%        \resumeItemListStart
%      \resumeItem{Developed path planning algorithm for robot navigation using B - Spline curves.}
%      \resumeItem{Optimized the algorithm by changing the parameters using RoS \& Gazebo.}
%      \resumeItemListEnd
%         \vspace{-15pt}
%        \resumeProjectHeading
%       {\textbf{Reliable Data Transfer (RDT) over UDP with various protocols} $|$ \textit{Prof.}  Praveen Tamanna }{}
%       \resumeItemListStart
%         \resumeItem{Simulated file transfer between UDP sockets in sender \& receiver using \textbf{Stop and Wait} \& \textbf{Go back N}  protocols.}
%         \resumeItem{Compared the average throughput by varying the retransmission timeout \& window size.}
%          \resumeItemListEnd
%          \vspace{-15pt}
%          \resumeProjectHeading
%       {\textbf{Ray Striker 2D} $|$ \textit{Prof.} Saurabh Joshi }{}
%       \resumeItemListStart
%         \resumeItem{Designed a 2D game naming \textbf{Ray Stricker 2D} in cpp using OOPs concepts.}
%         \resumeItem{Gained some fundamental knowledge of writing text-based user interface in a terminal independent manner.
% }
%          \resumeItemListEnd
        %   \vspace{-15pt}
        %          \resumeProjectHeading 
        % {\textbf{Database Model for Scientific Research Papers} $|$ \textit{Prof.} Manahor Kaul} {}
        % \resumeItemListStart
        % \resumeItem{Designed a Relational Database from scratch for Scientifc Research Papers.}
        % \resumeItem{Populated Database with the entries from the Research Paper Datasets using PostgreSQL.}
        % \resumeItemListEnd
        %   \vspace{-15pt}
        % \resumeProjectHeading 
        % {\textbf{Web Development - Publication Management System }$|$ \textit{Prof.} Manish Singh} {}
        % \resumeItemListStart
        % \resumeItem{Developed a website for managing the publications submitted to the conferences.}
        % \resumeItem{Used \textbf{Python}, \textbf{SQL} for developing the website.}
        % \resumeItemListEnd
%           \vspace{-15pt}
%         \resumeProjectHeading 
%         {\textbf{GraphIT - Compiler for custom language} $|$ \textit{Prof.} Ramakrishna Upadrasta} {}
%         \resumeItemListStart
%         \resumeItem{Developed a compiler for new langauge GraphIT , which uses graph as basic inbuilt data type \& has syntax for graph operations.
%         \resumeItem{Used \textbf{Flex}, \textbf{Bison}, \textbf{LLVM backend} for developing the compiler.}}
%         \resumeItemListEnd
%          \vspace{-5pt}



 
          % \vspace{-20pt}
  
% \vspace{-15pt}

% 
%
%-----------PROGRAMMING SKILLS-----------
\section{Technical Skills}
 \begin{itemize}[leftmargin=0.15in, label={}]
    % \small{
    \item{
     \textbf{Languages}{:  C, C++, Python, JavaScript} \\
     \vspace{5pt}
     \textbf{Frameworks \& Tools }{:  PyTorch, Tensorflow, ONNX, Hugging Face, Docker, Git, LaTeX, Markdown } \\
     \vspace{5pt}
     \textbf{HPC Frameworks}{: CUDA, Triton, DeepSpeed, TensorRT, ONNX Runtime}
     % \textbf{Operating Systems }{: Linux, Windows} \\
     % \textbf{Familiar}{: Tensorrt, JavaScript}\\
    }
    % }
 \end{itemize}
 \vspace{-16pt}

\section{Achievements}
    % \resumeSubHeadingListStart
    \resumeItemListStart
    % \resumeItem{Secured AIR \textbf{1156} in Open Category \& AIR \textbf{91} in EWS Category in IIT JEE Advanced 2020.}
    \resumeItem{Rank \textbf{1156} in IIT JEE Advanced 2020 among \textbf{150,000} participants.}
    % \resumeItem{Secured AIR \textbf{493} in Open Category \& AIR \textbf{45} in EWS Category in JEE Main B Planning 2020.}
    \resumeItem{Rank \textbf{493} in JEE Main B Planning 2020 among \textbf{59,000} participants.}
    \resumeItem{\textbf{99.75} percentile in IIT JEE Main 2020 among \textbf{1.1 Million} participants.}
    \resumeItem{AP EAMCET 2020 Rank \textbf{135} among \textbf{156,000} participants.}
    \resumeItem{TS EAMCET 2020 Rank \textbf{307} among \textbf{143,000} paricipants.}
    % \resumeItem{Solved around 700+ CP problems over multiple platforms \href{https://leetcode.com/KRS27/}{\color{Blue}LeetCode} , \href{https://www.codechef.com/users/sravanthk27}{\color{Blue}CodeChef} , \href{https://codeforces.com/profile/mrsravanth}{\color{Blue}CodeForces} , \href{https://auth.geeksforgeeks.org/user/ramasravanthkodavanti}{\color{Blue}GeeksForGeeks}.}
    \resumeItem{Solved around 700+ CP problems over multiple platforms \href{https://leetcode.com/sravanthkod/}{\color{Blue}LeetCode} , \href{https://www.codechef.com/users/sravanthk27}{\color{Blue}CodeChef} , \href{https://codeforces.com/profile/mrsravanth}{\color{Blue}CodeForces} , \href{https://auth.geeksforgeeks.org/user/ramasravanthkodavanti}{\color{Blue}GeeksForGeeks}.}
    \resumeItemListEnd
    
%-----------INVOLVEMENT---------------
\section{Leadership / Extracurricular}
    % \resumeSubHeadingListStart
    %     \resumeSubheading{Fraternity}{Spring 2020 -- Present}{President}{University Name}
            \resumeItemListStart
                \resumeItem{Worked as a Core member for \textbf{Epoch : AI-ML club of IITH}.}
                \resumeItem{Worked as an \textbf{Internship} \& \textbf{Placement} Coordinator for Office of Carrer Service of IITH.}
                \resumeItem{Member of IITH \textbf{Chess} team , Represented IITH at \textbf{InterIIT }\& various competitions.}
            \resumeItemListEnd
        
    \resumeSubHeadingListEnd


\end{document}
